\chapter{Similarity Search and the Moment Index \emph{(outline)}}\label{ch:similarity}
\chapterlede{``Find me moments that sound like this one'' --- across every
song in the library, in milliseconds, along a chosen facet (whole mix, just
the drums, just the vocals).}

\begin{incode}
Status: outline. To be written against \texttt{moment\_index.py},
\texttt{moments.py}, \texttt{descriptors.py}, \texttt{interactions.py}, and
the \texttt{/api/similar} endpoint.
\end{incode}

Planned contents:
\begin{itemize}
  \item \textbf{Exact top-$k$ retrieval.} Cosine scores as one matrix--vector
  product; complexity and memory layout; when brute force beats an index
  (spoiler: at this library's scale, always).
  \item \textbf{Faceted similarity.} Per-stem embedding blocks
  (\texttt{emb\_drums}, \texttt{emb\_bass}, \ldots) as projections of the
  product space; facet weighting as a diagonal metric.
  \item \textbf{Interpretable descriptors.} The 18-dimensional
  descriptor vector (RMS, onset density, spectral centroid, active
  fraction per stem); z-scoring; mixing learned and hand-crafted metrics.
  \item \textbf{Diversity constraints.} The per-song result cap as a greedy
  approximation to maximal marginal relevance.
  \item \textbf{Approximate methods} (for the future large library):
  inverted-file indices, HNSW graphs, product quantization --- what each
  trades away.
\end{itemize}
